% Reviewer 2
\reviewer%
\label{rev:2}

\begin{generalcomment}
	In general, the work is well-written. However, I have the following concerns.
\end{generalcomment}
\begin{revresponse}[Thank you for your feedback.]
	We have carefully addressed all the issues item by item as follows.
\end{revresponse}

\begin{revcomment}[label={comment:work-not-good}]
	The work is not really good.
\end{revcomment}
\begin{revresponse}
: (123)
\end{revresponse}

\begin{revcomment}
	You forgot to cite a very important reference (where I am an author)!
\end{revcomment}
\begin{revresponse}
	We are aware that citations on Google Scholar are very important to you.
	Therefore, we added reference~\cite{ReviewerReference}.

Also check out our article~\cite{Besser2020}.

	\printpartbibliography{ReviewerReference,Besser2020}
\end{revresponse}

\begin{revcomment}
	The template you are using for responding to my comments is excellent.
\end{revcomment}
\begin{revresponse}
	We agree.

	One cool feature is that comments are automatically labeled with the following pattern: \verb|comment:<rev_num>.<comment_num>|, which can then be used in a \verb|\ref|~command.
	So this comment has number~\ref{comment:2.3}.
	(Note that this also automatically changes if the \texttt{revnumtocomment} class option is enabled.)

	To automatically add the label in front of the number, we can use \verb|\autoref|: \autoref{comment:2.3}.

	It is also possible to set an additional custom label by passing the \texttt{label} option to the comment environment, e.g., \autoref{comment:work-not-good} was mean.
\end{revresponse}

